% (c) Copytight: ALL RIGHT RESERVED
% Author: Cwzx136 at GitHub
\documentclass[11pt,a4paper]{amsart} % looks nicer then common document
\usepackage{graphicx} % Required for inserting images
\usepackage{amsmath,amssymb,amsthm,mathtools,booktabs,array,enumitem,float}
\usepackage{fancyhdr}
\usepackage{xpatch} % For fixing sections
\title{Complete Classification of Primitive Integer Solutions to Ramanujan--Sato Series under a Motivic Period-Torsor Hypothesis}
\author{ZiXian Chang-Wang}
\date{August 2026}
% definitions for theorem-type envs
\newtheorem{theorem}{Theorem}[section]
\newtheorem{lemma}[theorem]{Lemma}
\newtheorem{proposition}[theorem]{Proposition}
\newtheorem{corollary}[theorem]{Corollary}
\theoremstyle{definition}
\newtheorem{definition}[theorem]{Definition}
\newtheorem{remark}[theorem]{Remark}
% Short for long and frequently used formulas
\newcommand{\Q}{\mathbb Q}
\newcommand{\Qbar}{\overline{\mathbb Q}}
\newcommand{\C}{\mathbb C}
\newcommand{\Z}{\mathbb Z}
\newcommand{\Mot}{\mathrm{Mot}}
\newcommand{\Sym}{\operatorname{Sym}}
\newcommand{\Hom}{\operatorname{Hom}}
\newcommand{\End}{\operatorname{End}}
\newcommand{\Pic}{\operatorname{Pic}}
\newcommand{\MT}{\operatorname{MT}}
\newcommand{\dR}{\mathrm{dR}}
\newcommand{\Betti}{\mathrm B}
\newcommand{\Per}{\mathrm{Per}}
\newcommand{\ord}{\operatorname{ord}}
\newcommand{\im}{\operatorname{Im}}
% Keep the running head short.
\makeatletter
\xapptocmd{\@sect}{\csname #1mark\endcsname{#7}}{}{}
\makeatother
\pagestyle{fancy}
\fancyhf{} 
\fancyhead[LO]{\footnotesize\leftmark}
\fancyhead[RO]{\footnotesize\thepage}
\fancyhead[RE]{\footnotesize ZiXian Chang-Wang}
\fancyhead[LE]{\footnotesize\thepage}
\renewcommand{\headrulewidth}{0.0pt}
\setlength{\headheight}{12pt}

\begin{document}
\begin{abstract}
Under the assumption that the Tannakian category generated by \(\mathrm{Sym}^2H^1(E)\oplus \mathbb{Q}(-1)\) satisfies the Grothendieck–André period torsor conjecture, we prove that any non‑degenerate \(1/\pi\)‑series relation forces the corresponding elliptic curve to have complex multiplication. When the Hauptmodul value is the rational number \(1/C\), the ring class number of the CM order is at most 2; hence the integer parameter \(C\) is restricted to finitely many values, and for a fixed \(C\), the primitive slope \((A:B)\) is unique. On this basis, combined with the known table of rational Ramanujan–Sato parameters for levels 1--4 and the complete enumeration of rational CM uniformizers with class number \(\le 2\), we obtain a conditional complete classification: there are 36 convergent primitive symbol classes for the standard branches 1A--4A, and 5 for the Euler–companion branches 1B, 2B, 3B, 4C; under the global sign flip \((A,B,\alpha)\mapsto(-A,-B,-\alpha)\), the total number is 41. We also formulate a higher-level finiteness theorem for genus-zero modular quotients, isolating the hypotheses needed to extend the argument beyond levels 1--4.
\end{abstract}
\maketitle
\tableofcontents
\section{Introduction}

In 1914, S. Ramanujan discovered a series of formulas for calculating $1/\pi$.\cite{Ram} The most famous one among them is \[\frac1\pi = \frac{2\sqrt{2}}{9801}\sum_{m=0}^\infty \frac{(4m)!}{(m!)^4}\frac{1103+26390m}{396^{4m}}\tag{$\star$}\] which adds approximately 8 decimal places of accuracy per term. Later, T. Sato provides a complete definition of such series. 

For \(\ell=1,2,3,4\), let
\[
\begin{aligned}
T_1(m)&=\binom{6m}{3m}\binom{3m}{m}\binom{2m}{m}
      =\frac{(6m)!}{(3m)!(m!)^3},\\
T_2(m)&=\binom{4m}{2m}\binom{2m}{m}^{2}
      =\frac{(4m)!}{(m!)^4},\\
T_3(m)&=\binom{3m}{m}\binom{2m}{m}^{2},\\
T_4(m)&=\binom{2m}{m}^{3}.
\end{aligned}
\]
Define
\[
{\mathfrak F}_\ell(X)=\sum_{m\ge0}T_\ell(m)X^m,\qquad
\Theta=X\frac{d}{dX},
\]
and
\[
{\mathfrak S}_\ell(A,B,C)=
\sum_{m\ge0}T_\ell(m)\frac{A+Bm}{C^m},
\qquad
A,B,C\in\Z,\quad C\ne0,\quad \gcd(A,B)=1.
\]
Infinite series of the above kind were called the standard branch level-$\ell$ Ramanujan-Sato series. The companion branch was given by Euler transformations. A simple observation leads to its hypergeometric representation: Within the radius of convergence,
\[
{\mathfrak S}_\ell(A,B,C)
=A {\mathfrak F}_\ell(C^{-1})+B\Theta {\mathfrak F}_\ell(C^{-1}).
\]

Until 2026, mathematicians had found 36 integer solutions to the level 1--4 standard branch. It's long open to show that they are exactly 36 primitive solutions.\cite{CCY} The goal of this work is to provide a conditional proof for the above conjecture. Our proof relies on the below hypothesis:

\subsection*{Period torsor hypothesis}

Let \(E/\Qbar\) be an elliptic curve,
\[
V=H^1(E),\qquad M_E=\Sym^2V\oplus\Q(-1).
\]
Denote by \(\mathfrak T_E=\langle M_E\rangle^\otimes\) the pure André motive Tannakian category generated by \(M_E\), and by \(\Omega^{\mathrm{mot}}_E\) its motivic period torsor. \cite{BCC}\cite{Andre}

\medskip
\noindent
\textbf{\((\mathrm{GPC}_{\mathrm{tor}})\).}
The complex point given by the comparison isomorphism is Zariski dense in \(\Omega^{\mathrm{mot}}_E\).
Equivalently, any polynomial period relation with coefficients in \(\Qbar\) that holds at the comparison isomorphism holds identically on the whole motivic period torsor.

\medskip

We use this density statement as an explicit motivic period-torsor hypothesis.  It is stronger in form than a fixed-codimension cycle statement such as $\mathrm{GPC}^k(X)$, and it is the form needed below to transport a numerical linear relation at the comparison point to an identity of matrix coefficients on the entire torsor.  The torsor-of-periods viewpoint and its relation with Grothendieck's period conjecture are discussed by Bost--Charles~\cite{BCC}.

\begin{theorem}[Main theorem]\label{thm:main}
Fix \(\ell\in\{1,2,3,4\}\).
Assume that \((\mathrm{GPC}_{\mathrm{tor}})\) holds for every elliptic curve \(E/\Qbar\) that appears in the proof.
Then there are only finitely many primitive integer triples \((A,B,C)\) satisfying
\[
{\mathfrak S}_\ell(A,B,C)=\frac{\alpha}{\pi},
\qquad \alpha\in\Qbar,
\]
and such that the series converges.
For each fixed \(C\), the primitive integer pair \((A,B)\) has at most one projective slope; hence at most one pair \((A,B)\) up to sign \((A,B)\) and \((-A,-B)\).
This statement includes the case \(\alpha=0\).
The proof does not use any upper bound on \([\Q(\alpha):\Q]\), so adding an arbitrary fixed degree condition \([\Q(\alpha):\Q]\le d\) does not change the conclusion.
\end{theorem}
In addition, in the standard branches and the Euler-companion branches defined below, we give the explicit level $1$--$4$ tables.  For example, the standard level $2A$ branch has the following twelve primitive representatives.  The positive-nome cases are

    \[\begin{array}{c|r|r|r|c}
N&C&A&B&\alpha\\ \hline
2&648&1&7&\dfrac92\\[1mm]
3&2304&1&8&2\sqrt3\\[1mm]
5&20736&1&10&\dfrac{9\sqrt2}{4}\\[1mm]
9&614656&3&40&\dfrac{49\sqrt3}{9}\\[1mm]
11&2509056&19&280&18\sqrt{11}\\[1mm]
29&24591257856&1103&26390&
\dfrac{9801\sqrt2}{4}
\end{array}\]
    and the negative-nome cases are
    \[\begin{array}{c|r|r|r|c}
N&C&A&B&\alpha\\ \hline
5&-1024&3&20&8\\
7&-3969&8&65&9\sqrt7\\
9&-12288&3&28&\dfrac{16\sqrt3}{3}\\[1mm]
13&-82944&23&260&72\\
25&-6635520&41&644&\dfrac{288\sqrt5}{5}\\[1mm]
37&-199148544&1123&21460&3528
\end{array}\]
The row $N=29$ gives Ramanujan's famous formula. \cite{Ram} ($\star$)

\section{Modular Parameterization}

\begin{lemma}[Hypergeometric representation]\label{lem:hyp}
    Let $(a_1,a_2,a_3,a_4)=\left(\frac16,\frac14,\frac13,\frac12\right)$ and $(\kappa_1,\kappa_2,\kappa_3,\kappa_4)=(1728,256,108,64)$. Then \[{\mathfrak F}_\ell(X)={}_3F_2\!\left(a_\ell,\frac12,1-a_\ell;1,1;\kappa_\ell X\right).\]
\end{lemma}
\begin{proof}
We give the case $\ell=2$; the other three cases follow from the corresponding multiplication formulas.  Gauss's multiplication formula with $n=4$ gives
\[ \Gamma (4 u) = \frac{4^{4 u - \frac{1}{2}}}{(2 \pi)^{3 / 2}} \text{ }
   \Gamma (u) \Gamma  \left( u + \frac{1}{4} \right) \Gamma  \left( u + \frac{1}{2} \right) \Gamma \left( u + \frac{3}{4}
   \right) . \]
Substituting $u=m+\frac14$ gives
\[
\Gamma(4m+1)=\frac{4^{4m+1/2}}{(2\pi)^{3/2}}
\Gamma\left(m+\frac14\right)
\Gamma\left(m+\frac12\right)
\Gamma\left(m+\frac34\right)
\Gamma(m+1).
\]
Since $\Gamma (4 m + 1) = (4 m) !$ and $\Gamma (m + 1) = m!$, while noticing
that $\Gamma (m + \alpha) / \Gamma (\alpha) = (\alpha)_m$, we obtain
\[ (4 m) ! = \frac{4^{4 m + \frac{1}{2}}}{(2 \pi)^{3 / 2}} \text{ } \Gamma
    \left( \frac{1}{4} \right) \Gamma  \left( \frac{1}{2}
   \right) \Gamma  \left( \frac{3}{4} \right) (1)_m  \left(
   \frac{1}{4} \right)_m \left( \frac{1}{2} \right)_m \left( \frac{3}{4}
   \right)_m . \]
Setting $m=0$ shows that the constant $\frac{4^{\frac{1}{2}}}{(2 \pi)^{3 / 2}}
\Gamma \left( \frac{1}{4} \right) \Gamma \left( \frac{1}{2} \right) \Gamma
\left( \frac{3}{4} \right) = 1$, and therefore
\[ T_2(m) = 4^{4 m} \frac{\left( \frac{1}{4}
   \right)_m \left( \frac{1}{2} \right)_m \left( \frac{3}{4}
   \right)_m}{(1)_m^3} . \] The remaining cases follow from the sextuple, triple, and double multiplication formulas, respectively.
\end{proof}

\begin{lemma}\label{lem:conv}
    For each \(\ell\in\{1,2,3,4\}\), there exists a constant \(c_\ell>0\) such that \[T_\ell(m)=c_\ell\frac{\kappa_\ell^m}{m^{3/2}}\left(1+O\!\left(\frac1m\right)\right).\] Therefore, the convergence of the series \({\mathfrak S}_\ell(A,B,C)\) is as follows: \[\begin{array}{c|c}
\text{Condition on }C & \text{Convergence of }{\mathfrak S}_\ell(A,B,C)\\ \hline
|C|>\kappa_\ell & \text{absolute convergence},\\
C=\kappa_\ell & \text{convergent iff }B=0,\\
C=-\kappa_\ell & \text{always convergent; conditionally if }B\ne0,\\
|C|<\kappa_\ell & \text{divergent}.
\end{array}\]
\end{lemma}
\begin{proof}
By Lemma~\ref{lem:hyp},
\[
T_\ell(m)=\kappa_\ell^m\frac{\Gamma(m+a_\ell)\Gamma(m+\frac12)\Gamma(m+1-a_\ell)}{\Gamma(a_\ell)\Gamma(\frac12)\Gamma(1-a_\ell)\Gamma(m+1)^3}.
\]
Stirling's formula, $m!\sim\sqrt{2\pi m}(m/e)^m$, gives
\[
\frac{\Gamma(m+r)}{\Gamma(m+1)}=m^{r-1}\left(1+O(m^{-1})\right).
\]
Hence the total exponent of $m$ is $(a_\ell-1)+(\frac12-1)-a_\ell=-\frac32$, and
\[
c_\ell=\frac1{\Gamma(a_\ell)\Gamma(\frac12)\Gamma(1-a_\ell)}>0.
\]

When \(|C|>\kappa_\ell\), the ratio test gives absolute convergence; when
\(|C|<\kappa_\ell\), the general term does not tend to zero.
When \(C=\kappa_\ell\), if \(B\ne0\), the terms have constant sign and are of order \(Bm^{-1/2}\), hence diverge; if \(B=0\), they are of order \(m^{-3/2}\), hence converge.
When \(C=-\kappa_\ell\), the terms alternate; if \(B\ne0\), their absolute values are of order \(m^{-1/2}\), giving conditional convergence; if \(B=0\), the convergence is absolute.
\end{proof}

We introduce the notation $q=e^{2\pi i\sigma}, \sigma\in\mathfrak H$. For \(\ell=1,2,3,4\), let $Y_\ell=X_0(\ell)^+$ be the Fricke quotient; for \(\ell=1\) we set \(Y_1=X(1)\). Campbell--Cooper--Ye take a Hauptmodul \(X_\ell=q+O(q^2)\) defined over \(\Q\), and a weight 2 modular form \(Z_\ell=1+O(q)\) satisfying
\[
Z_\ell
=\frac{1}{\sqrt{G_\ell(X_\ell)}}
q\frac{d}{dq}\log X_\ell,
\]
where
\[
\begin{array}{c|c|c}
\ell&G_\ell(X)&H_\ell(X)\\ \hline
1&1-1728X&240X\\
2&1-256X&48X\\
3&1-108X&24X\\
4&1-64X&16X.
\end{array}
\]

\begin{lemma}[Modular form equals hypergeometric solution]\label{lem:ZF}
    In a neighbourhood of \(q=0\), $Z_\ell(q)={\mathfrak F}_\ell(X_\ell(q))$. Consequently, \[q\frac{dX_\ell}{dq}=X_\ell Z_\ell\sqrt{1-\kappa_\ell X_\ell}.\]
\end{lemma}

\begin{proof}
Write
\[
Z_\ell=\sum_{n\ge0}U_\ell(n)X_\ell^n,\qquad U_\ell(0)=1.
\]
The recurrence of Campbell--Cooper--Ye, with \(G=1-\kappa X,\ H=hX\), is given by
\[
(n+1)^3U(n+1)
=\left(n+\frac12\right)
\bigl(h+\kappa n(n+1)\bigr)U(n).
\]
where $h_\ell=\kappa_\ell a_\ell(1-a_\ell)$. Therefore $h_\ell+\kappa_\ell n(n+1)=\kappa_\ell(n+a_\ell)(n+1-a_\ell)$, so
\[
\frac{U_\ell(n+1)}{U_\ell(n)}
=\kappa_\ell
\frac{(n+a_\ell)(n+\frac12)(n+1-a_\ell)}
{(n+1)^3}.
\]
This is exactly the ratio of consecutive terms of \(T_\ell(n)\) from Lemma \ref{lem:hyp}; since both have initial value \(1\), we get \(U_\ell(n)=T_\ell(n)\).

The second formula follows directly from the definition.
\end{proof}

\begin{lemma}[Algebraic parameters give algebraic elliptic curves]\label{lem:algE}
Let \(x\in\Qbar\) be a non-cuspidal value on \(Y_\ell\).
If \(X_\ell(\sigma)=x\), then one can choose a point
\[
\widetilde P\in X_0(\ell)(\Qbar)
\]
lying above the corresponding point; it parametrises an elliptic curve \(E_x/\Qbar\) together with a cyclic \(\ell\)-isogeny structure.
Specifically, if \(x\in\Q\), one can choose \(\widetilde P\) such that \([\Q(\widetilde P):\Q]\le2\)(with bound $1$ when \(\ell=1\)).
\end{lemma}
\begin{proof}
\(X_\ell\) is a coordinate on \(Y_\ell\simeq\mathbf P^1\) defined over \(\Q\),
thus an algebraic number \(x\) gives an algebraic point \(P_x\in Y_\ell\).
The natural map
\[
\pi_\ell:X_0(\ell)\longrightarrow X_0(\ell)^+
\]
is the quotient by the Fricke involution for \(\ell>1\), of degree \(2\); for \(\ell=1\) it is the identity.
The inverse image of an algebraic point under a finite map consists of algebraic points.
A geometric non-cuspidal point is represented, in moduli terms, by an elliptic curve together with a cyclic subgroup; to avoid automorphism issues of coarse moduli spaces, one may add an auxiliary full level \(N\ge3\), take the universal elliptic curve on a finite étale cover, and finally forget the auxiliary level.
When \(x\in\Q\), the residual field degree under the finite map is at most the degree of the map, giving the stated degree bound.
\end{proof}

\section{Period Torsor and Complex Multiplication}

Let $U_\ell=\mathbb P^1\setminus\{0,\kappa_\ell^{-1},\infty\}$. By Clausen's identity, ${\mathfrak F}_\ell(x)=G_\ell(x)^2$, where
\[
G_\ell(x)=
{}_2F_1\!\left(
\begin{matrix}
a_\ell/2,\ (1-a_\ell)/2\\
1
\end{matrix};\,\kappa_\ell x
\right).
\]
The function \(G_\ell\) satisfies the second-order equation
\[
\left[
\Theta^2-\kappa_\ell x
\left(\Theta+\frac{a_\ell}{2}\right)
\left(\Theta+\frac{1-a_\ell}{2}\right)
\right]\mathfrak G=0,
\]
while \({\mathfrak F}_\ell\)'s third-order equation is its symmetric square. For \(\ell=1,2,3,4\), this second-order local system is the elliptic Picard--Fuchs system for the corresponding arithmetic triangle group; on a fixed finite algebraic cover, it equals \(R^1\pi_*\Q\) for some family of elliptic curves.\\

\noindent\textbf{Notations. }We fix a finite algebraic cover \(U\) of the modular curve \(X_0(\ell)\) (or \(X(1)\) when \(\ell=1\)) on which the universal elliptic curve exists as a smooth proper morphism \(\pi:\mathcal E\to U\); for a point \(u\in U\), the fiber \(\mathcal E_u\) is the elliptic curve \(E_x\) corresponding to the modular parameter \(x\). The sheaf \(R^1\pi_*\Q\) denotes the first higher direct image of the constant sheaf \(\Q\) on \(\mathcal E\) in the complex analytic topology, whose fibre at \(u\) is the Betti cohomology \(H^1(\mathcal E_u(\C),\Q)\). Likewise, \(R^1\pi_*\Omega_{\mathcal E/U}^\bullet\) is the first derived direct image of the relative algebraic de Rham complex, a vector bundle on \(U\) whose fibre is the algebraic de Rham cohomology \(H^1_{\dR}(\mathcal E_u/\kappa(u))\), equipped with the Gauss–Manin connection. The comparison isomorphism \(\operatorname{comp}\) is the Grothendieck period isomorphism, which at each fibre gives a \(\C\)-linear identification \(H^1_{\dR}(E_x/\Qbar)\otimes_{\Qbar}\C \simeq H^1(E_x(\C),\C)\), and extends naturally to symmetric powers and tensor products. For a de Rham class \(s\in \Sym^2 H^1_{\dR}(E_x/\Qbar)\) and a Betti homology class \(\delta=\gamma^{\otimes 2}\in \Sym^2 H_1(E_x(\C),\Q)\), the pairing $\langle \operatorname{comp}(s),\delta\rangle$ is the natural tensor pairing between $\Sym^2H^1$ and $\Sym^2H_1$ (equivalently, via K\"unneth, the corresponding product-period pairing on $E_x\times E_x$); explicitly, if \(s=\omega_1\otimes\omega_2\), then \[\langle \operatorname{comp}(s),\gamma^{\otimes2}\rangle = \left(\int_\gamma \omega_1\right)\left(\int_\gamma \omega_2\right).\] Finally, \(F^2\Sym^2 H^1_{\dR}(E_x)\) denotes the second step of the Hodge filtration on the symmetric square; since \(F^1H^1_{\dR}(E_x)=H^{1,0}(E_x)\) is the space of holomorphic differentials, this subspace is \(\Sym^2(F^1H^1_{\dR}(E_x))\), consisting of classes of pure type \((2,0)\). \\

\begin{lemma}\label{lem:period}
Let $x\in\Qbar\setminus\{0,\kappa_\ell^{-1}\}$, and choose \(E_x\) as in Lemma \ref{lem:algE}. There exist two elements
\[
s_0,s_1\in
\Sym^2H^1_{\dR}(E_x/\Qbar),
\qquad
0\ne\delta\in
\Sym^2H_1(E_x(\C),\Q),
\]
such that
\begin{align*}
    \langle\operatorname{comp}(s_0),\delta\rangle
&=(2\pi i)^2{\mathfrak F}_\ell(x),\\
\langle\operatorname{comp}(s_1),\delta\rangle
&=(2\pi i)^2\Theta {\mathfrak F}_\ell(x).
\end{align*}
Moreover, on the real convergence interval $-\kappa_\ell^{-1}\le x<\kappa_\ell^{-1},x\ne0,$ the vectors \(s_0,s_1\) are linearly independent. One may also take $$s_0\in F^2\Sym^2H^1_{\dR}(E_x)\setminus \{0\}.$$
\end{lemma}

\begin{proof}
In the modular curve with full auxiliary level, denote the universal elliptic curve by \(\pi:\mathcal E\to U\), and set
\[
\mathcal H_{\dR}=R^1\pi_*\Omega^\bullet_{\mathcal E/U},
\qquad
\mathcal L=F^1\mathcal H_{\dR}.
\]
The rational \(q\)-expansion of the weight 2 modular form \(Z_\ell\), by the \(q\)-expansion principle, gives an algebraic section defined over \(\Q\): $\mathbf z\in H^0(U,\mathcal L^{\otimes2})\subset H^0(U,\Sym^2\mathcal H_{\dR})$.
On the Tate curve the invariant differential is $dt/t=2\pi i\,dz$.  With the standard normalization in which $Z_\ell=1+O(q)$ and a local flat cycle $\gamma$ corresponding to one turn around the Tate parameter, the section $\mathbf z=Z_\ell(dt/t)^{\otimes2}$ therefore satisfies
$\langle\operatorname{comp}(\mathbf z),\gamma^{\otimes2}\rangle=(2\pi i)^2Z_\ell$. At \(X_\ell=x\), set
\[
D=X_\ell\frac{\partial}{\partial X_\ell}.
\]
Because the Gauss--Manin connection commutes with period integration, $D\langle\operatorname{comp}(\mathbf z),\gamma^{\otimes2}\rangle=\left\langle\operatorname{comp}(\nabla_D\mathbf z),\gamma^{\otimes2}\right\rangle$, since \(\gamma\) is flat.
By Lemma \ref{lem:ZF}, $D Z_\ell=\Theta {\mathfrak F}_\ell(X_\ell)$. Choose
\[
s_0=\mathbf z_x,\qquad
s_1=(\nabla_D\mathbf z)_x,\qquad
\delta=\gamma_x^{\otimes2},
\]
gives the two equalities at the point $x$.\\
We now prove linear independence. Locally write $\mathbf z=f\,\omega^{\otimes2}$ where $\omega$ generates $\mathcal L$. On the real convergence interval, Clausen's formula gives ${\mathfrak F}_\ell(x)=G_\ell(x)^2>0$ so \(f(x)\ne0\).\\
The Kodaira--Spencer map $\rho: TU \longrightarrow R^1\pi_*T_{\mathcal E/U}$ of the modular curve is an isomorphism at nondegenerate points; hence $\nabla_D\omega\not\equiv0\pmod{\mathcal L}$ as long as \(D\ne0\). Modulo \(F^2\Sym^2\mathcal H_{\dR}= \Sym^2(F^1 H^1_{\mathrm{dR}}(E_x))\) we have $\nabla_D\mathbf z\equiv2f\,\omega\otimes\nabla_D\omega\ne0$. Since \(\mathbf z\in F^2\), \(\mathbf z_x\) and \((\nabla_D\mathbf z)_x\) are linearly independent. Finally \(s_0\in\mathcal L_x^{\otimes2}=F^2\Sym^2H^1_{\dR}(E_x)\).
\end{proof}

\begin{lemma}\label{lem:coeff}
Let \(\mathfrak T\) be a semisimple neutral Tannakian category over an algebraically closed field of characteristic zero, and let \(G\) be its Tannaka group.
For a simple object \(W\), the matrix coefficient map
\[
W_{\dR}\otimes W_{\Betti}^{\vee}
\longrightarrow \mathcal O(G)
\]
is injective.
If \(W_1,W_2\) are non-isomorphic simple objects, then their matrix coefficient subspaces in \(\mathcal O(G)\) intersect trivially.
The same conclusion holds for a \(G\)-torsor: choose any geometric base point to identify the torsor with \(G\).
\end{lemma}

\begin{proof}
For a simple representation \(W\), the image of the group algebra in \(\End(W)\) generates all of \(\End(W)\); this is the finite-dimensional form of Burnside's density theorem.
If a tensor \(u\otimes\lambda\) gives a matrix coefficient \(\lambda(gu)\) that vanishes for all \(g\in G\), then \(\lambda(Au)=0\) for all \(A\in\End(W)\), hence \(u=0\) or \(\lambda=0\), and hence the matrix coefficient map is injective. The matrix coefficient spaces of different simple representations are different simple subcoalgebras of the representative function coalgebra. In a semisimple coalgebra, distinct simple subcoalgebras intersect trivially.
\end{proof}

\begin{lemma}\label{lem:noncm}
If \(E/\Qbar\) has no complex multiplication, let \(V=H^1(E)\). Then $M=\Sym^2V$
is a simple object, and $\Hom_{\Mot}(\Q(-1),M)=0$.
\end{lemma}

\begin{proof}
No CM is equivalent to \(\End_{\mathrm{HS}}(V)=\Q\). In that case the derived group of the Hodge group acts as \(\mathrm{SL}_2\) on \(V\), and the three-dimensional representation \(\Sym^2\) is irreducible, so \(\Sym^2V\) has no nontrivial Hodge substructure; a motivic subobject would give a Hodge substructure, hence \(M\) is simple. On the other hand,
\[
V\otimes V=\Sym^2V\oplus\bigwedge^2V,
\qquad
\bigwedge^2V\simeq\Q(-1).
\]
Since \(\End_{\mathrm{HS}}(V)=\Q\), the Tate line in \(V\otimes V\) is unique and given by the polarisation; it lies in the alternating square \(\bigwedge^2V\), not in \(\Sym^2V\). Hence the stated $\Hom$ space is zero.
\end{proof}

\begin{theorem}[\(1/\pi\) relation forces CM]\label{thm:forceCM}
Let $x$ lie in the real convergence interval
\[
-\kappa_\ell^{-1}\le x<\kappa_\ell^{-1},\qquad x\ne0,
\]
and suppose
\[
u{\mathfrak F}_\ell(x)+v\Theta {\mathfrak F}_\ell(x)=\frac{\beta}{\pi},
\qquad
u,v,\beta\in\Qbar,\quad (u,v)\ne(0,0).
\]
Then under \((\mathrm{GPC}_{\mathrm{tor}})\), the corresponding elliptic curve \(E_x\) must have complex multiplication.
This conclusion includes the case when \(\beta=0\).
\end{theorem}

\begin{proof}
Take \(s_0,s_1,\delta\) from Lemma \ref{lem:period}. Multiply the equality by \((2\pi i)^2\), and use \(1/\pi=2i/(2\pi i)\), obtaining $\left\langle\operatorname{comp}(us_0+vs_1),\delta\right\rangle-2i\beta(2\pi i)=0$. The left-hand side is the sum of two matrix coefficients that come from $M=\Sym^2H^1(E_x)$ and $T=\Q(-1)$. By $(\mathrm{GPC}_{\mathrm{tor}})$ the comparison point is Zariski dense in the motivic period torsor.  Therefore this regular linear function, which vanishes at the comparison point, vanishes identically on the torsor.\\
Assume, for contradiction, that \(E_x\) is non-CM. By Lemma \ref{lem:noncm}, \(M\) and \(T\) are non-isomorphic simple objects. By Lemma \ref{lem:coeff}, their matrix coefficient spaces are direct summands and separate, so the two components must vanish separately, i.e. $2i\beta=0,(us_0+vs_1)\otimes\delta=0$. Since \(\delta\ne0\) and \(s_0,s_1\) are linearly independent, the second equation forces \(u=v=0\), a contradiction.
\end{proof}

\begin{theorem}[Relation space at CM points is at most one-dimensional]\label{thm:unique}
Let $x$ lie in the real convergence interval of Lemma~\ref{lem:period}, and assume that $E_x$ has CM.  Set
\[
p_0=(2\pi i)^2{\mathfrak F}_\ell(x),\qquad
p_1=(2\pi i)^2\Theta {\mathfrak F}_\ell(x),\qquad
t=2\pi i.
\]
Under \((\mathrm{GPC}_{\mathrm{tor}})\), the \(\Qbar\)-linear relation space among the triple \((p_0,p_1,t)\) is at most one-dimensional.
\end{theorem}

\begin{proof}
Extend the coefficient field to an algebraically closed field containing the CM field \(K\). If $V=H^1(E_x)=L\oplus\overline L$ then $M=\Sym^2V
=L^{\otimes2}\oplus(L\otimes\overline L)
\oplus\overline L^{\otimes2}$, the polarization gives $L\otimes\overline L\simeq T=\Q(-1)$. The other two one-dimensional objects have Hodge types \((2,0)\) and \((0,2)\), respectively, and are non-isomorphic to the Tate object.\\
The Betti tensor from Lemma \ref{lem:period} may be taken as \(\delta=\gamma^{\otimes2}\), where \(0\ne\gamma\in H_1(E_x,\Q)\). Under the canonical decomposition \(V_{\Betti}=L_{\Betti}\oplus\overline L_{\Betti}\),
the projections of \(\gamma\) to both components are nonzero; otherwise complex conjugation exchanges the two CM characters and fixes the rational vector \(\gamma\), forcing \(\gamma=0\).
Thus the projections of $\delta$ to $L^{\otimes2}$ and $\overline L^{\otimes2}$ are also nonzero.  Moreover, the mixed projection to
$L\otimes\overline L\simeq T$ is the product of the two nonzero projections of $\gamma$, and is therefore nonzero.  Consequently a Tate de Rham vector has a nonzero Tate matrix coefficient against $\delta$.\\
Suppose $ap_0+bp_1+ct=0.$ By GPC and Lemma \ref{lem:coeff}, the two non-Tate matrix coefficients cannot be cancelled by the Tate matrix coefficient. Therefore $as_0+bs_1\in T_{\dR}\subset M_{\dR}$. Conversely, whenever \(as_0+bs_1\in T_{\dR}\), its Tate matrix coefficient determines a unique \(c\).\\
Hence the relation space has the same dimension as $W\cap T_{\dR}$, where $W=\operatorname{span}_{\Qbar}\{s_0,s_1\}$. By Lemma \ref{lem:period}, \(\dim W=2\). Also $0\ne s_0\in F^2M_{\dR}$, while the Tate object \(T=\Q(-1)\) has Hodge type \((1,1)\), so $F^2T_{\dR}=0$.\\
Thus \(s_0\notin T_{\dR}\), so $\dim(W\cap T_{\dR})\le1$; hence the relation space is at most one-dimensional.
\end{proof}

\begin{corollary}[Unique primitive slope for fixed parameter]\label{cor:slope}
Fix $x=1/C$ with
$-\kappa_\ell^{-1}\le x<\kappa_\ell^{-1}$ and $x\ne0$.  If there are two relations
\[
A_i{\mathfrak F}_\ell(x)+B_i\Theta {\mathfrak F}_\ell(x)
=\frac{\alpha_i}{\pi},
\qquad i=1,2,
\]
then
\[
A_1B_2-A_2B_1=\det\left(\begin{matrix}
    A_1&B_1\\A_2&B_2
\end{matrix}\right)=0.
\]
If both pairs of integers are primitive, then $(A_2,B_2)=\pm(A_1,B_1)$.
\end{corollary}

\begin{proof}
This follows from Theorem~\ref{thm:unique}: If the determinant were nonzero, the two relations would be linearly independent among the three numbers \({\mathfrak F}_\ell(x),\Theta {\mathfrak F}_\ell(x),1/\pi\), contradicting Theorem \ref{thm:unique}. In addition, primitive integer vectors representing the same rational projective point can differ only by a sign.
\end{proof}

\section{Rational Hauptmodul values and ring class numbers}

\begin{lemma}\label{lem:jdegree}
Let $x=X_\ell(\sigma)\in\Q$ and \(\sigma\) be a CM point.
Then $[\Q(j(\sigma)):\Q]\le2$. Moreover, when \(\ell=1\), the bound is \(1\).
\end{lemma}

\begin{proof}
\(x\in\Q\) gives a rational point on \(Y_\ell\).
By Lemma \ref{lem:algE}, one can choose a preimage \(\widetilde P\) on \(X_0(\ell)\) such that $[\Q(\widetilde P):\Q]\le2$ and the bound is \(1\) when \(\ell=1\). The function \(j\) is defined on \(X_0(\ell)\) and defined over \(\Q\), hence $j(\sigma)=j(\widetilde P)\in\Q(\widetilde P)$.
\end{proof}

\subsection*{Ring-class reciprocity}
Let $\mathcal O_D$ be an imaginary quadratic order of discriminant $D<0$.  Proper invertible $\mathcal O_D$-ideal classes act simply transitively on the CM elliptic curves with endomorphism ring $\mathcal O_D$.  Equivalently, the ring class polynomial
\[
H_D(Y)=\prod_{[\mathfrak a]\in\Pic(\mathcal O_D)}\bigl(Y-j(\mathfrak a)\bigr)
\]
is the minimal polynomial over $\Q$ of any one of its roots and has degree $h(D)=\#\Pic(\mathcal O_D)$; see, for example, \cite{Cox}.  This is the ring-class version of the usual complex multiplication reciprocity law and is the form used below.

\begin{lemma}\label{lem:h2}
Let \(\mathcal O_D\) be the imaginary quadratic order corresponding to \(\sigma\). Then $h(D)=\#\Pic(\mathcal O_D)=[\Q(j(\sigma)):\Q]\le2.$
\end{lemma}

\begin{proof}
By ring-class reciprocity, $H_D(Y)\in\mathbf Z[Y]$ is irreducible over $\Q$ and has degree $h(D)$.  Therefore
$[\Q(j(\sigma)):\Q]=h(D)$.  Lemma~\ref{lem:jdegree} now gives $h(D)\le2$.
\end{proof}

\begin{lemma}\label{lem:phi}
As $f\to\infty$, one has $\varphi(f)\to\infty$.
\end{lemma}

\begin{proof}
Fix $H>0$.  If $p^a\mid f$, then
\[
\varphi(p^a)=p^{a-1}(p-1)\le \varphi(f).
\]
Thus $\varphi(f)\le H$ implies $p-1\le H$ for every prime divisor $p$ of $f$, and for each such $p$ it also bounds the exponent $a$.  Hence only finitely many integers $f$ satisfy $\varphi(f)\le H$.  This is equivalent to $\varphi(f)\to\infty$.
\end{proof}

\begin{lemma}[Bounded ring class numbers give finitely many orders]\label{lem:finiteD}
For any fixed \(H\), there are only finitely many negative discriminants \(D\) satisfying \(h(D)\le H\).
\end{lemma}

\begin{proof}
First assume \(D=D_K\) is a fundamental discriminant. The Siegel--Heilbronn lower bound says that for every \(\varepsilon>0\), there exists \(c(\varepsilon)>0\) such that
\[
h(D_K)\ge
c(\varepsilon)|D_K|^{1/2-\varepsilon}.
\]
Thus as the absolute value of a fundamental discriminant tends to infinity, the class number tends to infinity; there are only finitely many imaginary quadratic fields of bounded class number.\\
In general write $D=f^2D_K$. The class number formula for orders is
\[
h(f^2D_K)
=
\frac{h(D_K)f}
{[\mathcal O_K^\times:\mathcal O_f^\times]}
\prod_{p\mid f}
\left(1-\frac{(\frac{D_K}{p})}{p}\right).
\]
Since
\[
1-\frac{(\frac{D_K}{p})}{p}\ge1-\frac1p,
\]
we have
\[
f\prod_{p\mid f}
\left(1-\frac{(\frac{D_K}{p})}{p}\right)
\ge
f\prod_{p\mid f}\left(1-\frac1p\right)
=\varphi(f).
\]
The unit index is at most \(3\), while by Lemma \ref{lem:phi}, \(\varphi(f)\to\infty\). Thus for fixed \(D_K\), \(h(f^2D_K)\to\infty\) as \(f\to\infty\). Combining with the finite number of \(D_K\) from the previous paragraph, the result follows.
\end{proof}

\section{Main Theorem and the Companion Branch}

\begin{proof}[Proof of Theorem \ref{thm:main}]
By Lemma \ref{lem:conv}, convergence requires $|C|\ge\kappa_\ell$. First treat \(C=\kappa_\ell\). Convergence then forces \(B=0\), and primitivity gives \(A=\pm1\). Thus this boundary contributes at most one sign class.\\
Now assume \(C\ne\kappa_\ell\), and set $x=1/C\in\Q$. Then $-\kappa_\ell^{-1}\le x<\kappa_\ell^{-1}$ and $x\neq 0$. If
\[
A{\mathfrak F}_\ell(x)+B\Theta {\mathfrak F}_\ell(x)=\frac{\alpha}{\pi},
\]
then Theorem \ref{thm:forceCM} shows that the corresponding elliptic curve \(E_x\) has CM. Let its CM order discriminant be \(D\). Since \(x\in\Q\), Lemma \ref{lem:h2} gives $h(D)\le2$. By Lemma \ref{lem:finiteD}, there are only finitely many possible \(D\).\\
For each fixed \(D\), the isomorphism classes of CM elliptic curves of discriminant \(D\) are parametrised by the finite set \(\Pic(\mathcal O_D)\).
Adding a cyclic \(\ell\)-subgroup on \(X_0(\ell)\) gives only finitely many choices. Therefore the corresponding \(Y_\ell\)-points and their Hauptmodul values $x=X_\ell(\sigma)$ are finite in number.
Among them, those satisfying \(x=1/C\) with \(C\in\Z\) are certainly finite.
Hence the possible \(C\)'s are finite.\\
For each \(C\ne\kappa_\ell\), Corollary \ref{cor:slope} shows that the primitive integer pair \((A,B)\) has at most one sign class. Together with the already separately handled case \(C=\kappa_\ell\), we obtain finiteness of triples \((A,B,C)\).
Once \((A,B,C)\) is fixed,
\[
\alpha=
\pi\bigl(A{\mathfrak F}_\ell(C^{-1})+
B\Theta {\mathfrak F}_\ell(C^{-1})\bigr)
\]
is uniquely determined.
\end{proof}

\begin{corollary}[Finite degree parameter curve version]\label{cor:companion}
Let \(Y/\Q\) be a smooth projective algebraic curve equipped with two nonconstant rational functions defined over \(\Q\),
\[
y:Y\to\mathbf P^1,\qquad j:Y\to\mathbf P^1,
\]
and suppose \(\deg(y)=d<\infty\). Assume that for a companion branch, the parameter \(x\) comes from \(Y\), and its Picard--Fuchs system is still the symmetric square of an elliptic \(H^1\).
If an integer denominator condition gives $y(P)=1/C\in\Q$, then $[\Q(j(P)):\Q]\le d$. Under \((\mathrm{GPC}_{\mathrm{tor}})\), points satisfying the corresponding \(1/\pi\) relation must be CM points, so $h(D)\le d$, and therefore primitive integer solutions still form a finite set.
\end{corollary}

\begin{proof}
The rational number \(y(P)\) gives a rational point on \(\mathbb P^1\).
Since a nonconstant rational function on a smooth projective curve defines a finite morphism, the residue field degree of any closed point above a rational value of $y$ is at most $d$; hence $[\Q(P):\Q]\le d$. Since \(j\) is defined over \(\Q\), $j(P)\in\Q(P)$, so \([\Q(j(P)):\Q]\le d\). The rest of the argument is the same as in the main theorem: the period relation forces CM, CM theory identifies the \(j\)-degree with the ring class number, and bounded ring class numbers give finitely many discriminants and finitely many parameter points.
\end{proof}

\section{Higher-level finiteness mechanism}\label{sec:higher}

The motivic part of the preceding argument does not depend on the factorial closed forms specific to levels $1$--$4$.  The essential inputs are a modular parameter, a weight-two Hodge section, a finite map to a modular curve carrying an elliptic family, and control of the degree of a lift.

\begin{theorem}[Higher-level genus-zero version]\label{thm:higher}
Let $Y/\Q\simeq\mathbf P^1$ be a modular quotient with Hauptmodul $X=q+O(q^2)$ defined over $\Q$, and suppose that there is a finite modular curve $\widetilde Y\to Y$ of degree $d$ on which the modular $j$-function is defined; after a further fixed fine-level cover we may use a universal elliptic curve to construct the Gauss--Manin connection.  Let $Z=1+O(q)$ be a weight-two modular form defined over $\Q$ and write locally at the cusp
\[
Z(\tau)=F(X(\tau)),\qquad F(X)=\sum_{n\ge0}T(n)X^n,
\qquad \Theta=X\frac{d}{dX}.
\]
Fix a rational value $x=1/C$ for which the point is noncuspidal, $X$ is a local parameter on the chosen cover, $Z(x)\ne0$, and both $F(x)$ and $\Theta F(x)$ are represented by the convergent branch under consideration.  Assume the motivic period-torsor density hypothesis for the corresponding elliptic curve.

If
\[
uF(x)+v\Theta F(x)=\frac{\beta}{\pi},
\qquad u,v,\beta\in\Qbar,\quad (u,v)\ne(0,0),
\]
then the corresponding elliptic curve has complex multiplication.  Moreover its CM order has class number at most $d$.  Consequently only finitely many such rational parameters $x$ occur.  At every parameter for which the two de Rham vectors obtained from $Z$ and its Gauss--Manin derivative are independent, the projective slope $[u:v]$ is unique.
\end{theorem}

\begin{proof}
The section $Z$ determines an algebraic section of $F^2\Sym^2H^1_{\mathrm{dR}}$, and its Gauss--Manin derivative realizes $\Theta F$ exactly as in Lemma~\ref{lem:period}.  The non-CM matrix-coefficient argument of Theorem~\ref{thm:forceCM} therefore applies verbatim.  If $x\in\Q$, a point above it on $\widetilde Y$ has residue degree at most $d$, and hence so does its $j$-invariant.  For a CM point this degree equals the ring class number, so $h(D)\le d$.  Lemma~\ref{lem:finiteD} gives only finitely many possible orders, and each order yields only finitely many points on the fixed modular cover.  The uniqueness statement follows from the CM decomposition used in Theorem~\ref{thm:unique} whenever the two de Rham vectors are independent.
\end{proof}

\begin{remark}
For a Fricke quotient $Y=X_0(N)/\langle w_N\rangle$, one has $d\le2$ (and $d=1$ for $N=1$).  Campbell--Cooper--Ye treat all genus-zero groups $\Gamma_0(N)^+$ in their classification~\cite{CCY}.  Theorem~\ref{thm:higher} therefore extends the conditional \emph{finiteness and unique-slope mechanism} to those higher levels, provided the stated local nonvanishing, independence, and convergence hypotheses are verified.  It does not by itself supply complete higher-level integer tables: those require an explicit enumeration of the relevant CM Hauptmodul values and a separate convergence analysis for each branch.
\end{remark}

\section{Complete conditional classification}

Sections~1--5 prove the level $1$--$4$ finiteness theorem, while Section~\ref{sec:higher} records the higher-level extension. In this section we add the explicit enumeration of low-class-number CM points to the argument, thereby obtaining the primitive integer tables for the standard
\(1A,2A,3A,4A\) branches and their Euler-companion branches.
%7.1->ast, 7.2->clubsuit, 7.3->spadesuit
\subsection*{Standard A-branch}
Campbell--Cooper--Ye's Theorem~2.4 gives a uniform formula
\cite{CCY}: if both $X=1/C$ and the projective slope $\lambda=A/B$ are rational, then on the corresponding positive and negative nome branches,
\[
\sum_{n\ge0}T_\ell(n)(n+\lambda)X^n
=\frac{1}{2\pi\sqrt{1-\kappa_\ell X}}
\begin{cases}
\sqrt{\ell/N},&\text{positive nome},\\[1mm]
\sqrt{L_\ell/N},&\text{negative nome},
\end{cases}
L_\ell=\frac{4\ell}{\gcd(4,\ell)}.
\tag{$\ast$}
\]
For these fixed Clausen branches, Chen--Glebov--Goenka explain that the possible rational CM uniformizer values can be enumerated from the known imaginary quadratic orders of class number at most two, and their single-sum cases match the rational level $1$--$4$ tables~\cite[Sec.~10]{CGG}.  We use this explicit CM enumeration, together with the Campbell--Cooper--Ye tables, as the external arithmetic input needed to pass from finiteness to the stated complete list.

\begin{proposition}[From \((X,\lambda,N)\) to primitive integer quadruples]\label{prop:normalizeA}
Let \(X=1/C\in\mathbf Q\), \(C\in\mathbf Z\setminus\{0\}\), and write
\[
\lambda=\frac AB,\qquad A,B\in\mathbf Z,\quad B>0,
\quad \gcd(A,B)=1.
\]
If the series corresponding to $(\ast)$ converges, then
\[
\sum_{n\ge0}T_\ell(n)\frac{A+Bn}{C^n}=\frac{\alpha}{\pi},
\]
where
\[
\alpha=
\frac{B}{2\sqrt{1-\kappa_\ell/C}}
\begin{cases}
\sqrt{\ell/N},&\text{positive nome},\\[1mm]
\sqrt{L_\ell/N},&\text{negative nome}.
\end{cases}
\tag{$\clubsuit$}
\]
In particular, the \(\alpha\) in the tables automatically lies in an extension of degree at most two.
\end{proposition}

\begin{proof}
Multiply ($\ast$) by \(B\), and use \(B(n+A/B)=A+Bn\). Since \(C,N,\ell,L_\ell\in\mathbf Q\), the right-hand side contains only one square root, so \([\mathbf Q(\alpha):\mathbf Q]\le2\).
\end{proof}

\begin{lemma}\label{lem:Apositiveboundary}
Under \((\mathrm{GPC}_{\mathrm{tor}})\), the standard \(A\) branch at
\(C=\kappa_\ell\) yields no primitive integer \(1/\pi\) solutions.
\end{lemma}

\begin{proof}
By Lemma~\ref{lem:conv}, at \(C=\kappa_\ell\) convergence forces \(B=0\),
and primitivity further gives \(A=\pm1\). The point
\(x=\kappa_\ell^{-1}\) is a finite elliptic branch point of the modular curve. After passing to a finite cover that removes orbifold automorphisms, the weight-two Hodge section still gives a nonzero vector $0\ne s_0\in F^2\Sym^2H^1_{\mathrm{dR}}(E_x)$. 
If \({\mathfrak F}_\ell(x)=\alpha/\pi\), then the same period torsor argument as in Theorem~\ref{thm:forceCM} yields an identity between this matrix coefficient of \(\Sym^2H^1(E_x)\) and a Tate matrix coefficient.
In the CM decomposition, the non-Tate matrix coefficients must vanish separately, so \(s_0\) would have to lie in the Tate line
\(\mathbf Q(-1)_{\mathrm{dR}}\). However, $F^2\mathbf Q(-1)_{\mathrm{dR}}=0$, contradicting $s_0\ne0$. 
\end{proof}

\begin{theorem}[Complete primitive integer classification of standard \(A\) branches]\label{thm:Acomplete}
Assume \((\mathrm{GPC}_{\mathrm{tor}})\) holds for the elliptic curves appearing in the proof.
Up to the overall sign
\[
(A,B,\alpha)\sim(-A,-B,-\alpha),
\]
the convergent primitive integer solutions of the standard \(1A,2A,3A,4A\) branches are exactly the \(36\) entries listed in
Tables~\ref{tab:1A}--\ref{tab:4Acomplete}:
\[
11+12+9+4=36.
\]
Except for the Bauer formula \((A,B,C)=(1,4,-64)\) in \(4A\), all other \(35\) entries are absolutely convergent.
\end{theorem}

\begin{proof}
Take any convergent primitive solution. If \(C=\kappa_\ell\), it is excluded by Lemma
\ref{lem:Apositiveboundary}. Otherwise set \(x=1/C\).
Theorem~\ref{thm:forceCM} forces \(E_x\) to have CM; since \(x\in\mathbf Q\),
Lemma~\ref{lem:h2} gives \(h(D)\le2\).\\
Chen--Glebov--Goenka give a complete enumeration of rational uniformizers for these low-class-number imaginary quadratic orders
\cite[Sec.~10]{CGG}; Campbell--Cooper--Ye's tables for levels \(1\)--\(4\)
\cite[Tables~5--11]{CCY} give the corresponding rational parameters \(X\) and slopes \(\lambda\).
Applying simultaneously the conditions \(X=1/C\), \(C\in\mathbf Z\), \(\lambda\in\mathbf Q\), and the convergence conditions of Lemma~\ref{lem:conv}, yields exactly the following \(36\) rows.
There are two additional rational CM candidates
\[
(\ell;A,B,C)=(2;1,5,-144),(3;4,15,-27),
\tag{$\spadesuit$}
\]
but they satisfy respectively \(|C|<256\) and \(|C|<108\), so the original \(A\)-series diverges and they are not counted.\\
For each surviving \(C\), Corollary~\ref{cor:slope} says that the primitive slope is unique up to the overall sign;
and Proposition~\ref{prop:normalizeA} gives the nonzero \(\alpha\) in the table. Hence there are neither omissions nor repetitions.
Finally, the convergence type is read off directly from Lemma~\ref{lem:conv}; only
\(C=-\kappa_4=-64\) with \(B\ne0\), the Bauer case, is conditionally convergent.
\end{proof}

\begin{table}[htbp]
\centering
\caption{All primitive integer solutions for Level \(1A\) (taking \(B>0\))}\label{tab:1A}
\footnotesize
\setlength{\tabcolsep}{4pt}
\renewcommand{\arraystretch}{1.14}
\begin{tabular}{@{}c r r r r l@{}}
\toprule
nome & \(N\) & \(A\) & \(B\) & \(C\) & \(\alpha\)\\
\midrule
+&2&3&28&8000&\(5\sqrt5\)\\
+&3&1&11&54000&\(5\sqrt{15}/6\)\\
+&4&5&63&287496&\(11\sqrt{33}/4\)\\
+&7&8&133&16581375&\(85\sqrt{255}/54\)\\
-&7&8&63&-3375&\(5\sqrt{15}\)\\
-&11&15&154&-32768&\(32\sqrt2\)\\
-&19&25&342&-884736&\(32\sqrt6\)\\
-&27&31&506&-12288000&\(160\sqrt{30}/9\)\\
-&43&263&5418&-884736000&\(640\sqrt{15}/3\)\\
-&67&10177&261702&-147197952000&\(1760\sqrt{330}\)\\
-&163&13591409&545140134&-262537412640768000&\(426880\sqrt{10005}\)\\
\bottomrule
\end{tabular}
\end{table}

\begin{table}[htbp]
\centering
\caption{All primitive integer solutions for Level \(2A\) (taking \(B>0\))}\label{tab:2A}
\footnotesize
\setlength{\tabcolsep}{4pt}
\renewcommand{\arraystretch}{1.14}
\begin{tabular}{@{}c r r r r l@{}}
\toprule
nome & \(N\) & \(A\) & \(B\) & \(C\) & \(\alpha\)\\
\midrule
+&2&1&7&648&\(9/2\)\\
+&3&1&8&2304&\(2\sqrt3\)\\
+&5&1&10&20736&\(9\sqrt2/4\)\\
+&9&3&40&614656&\(49\sqrt3/9\)\\
+&11&19&280&2509056&\(18\sqrt{11}\)\\
+&29&1103&26390&24591257856&\(9801\sqrt2/4\)\\
-&5&3&20&-1024&\(8\)\\
-&7&8&65&-3969&\(9\sqrt7\)\\
-&9&3&28&-12288&\(16\sqrt3/3\)\\
-&13&23&260&-82944&\(72\)\\
-&25&41&644&-6635520&\(288\sqrt5/5\)\\
-&37&1123&21460&-199148544&\(3528\)\\
\bottomrule
\end{tabular}
\end{table}

\begin{table}[htbp]
\centering
\caption{All primitive integer solutions for Level \(3A\) (taking \(B>0\))}\label{tab:3A}
\footnotesize
\setlength{\tabcolsep}{4pt}
\renewcommand{\arraystretch}{1.14}
\begin{tabular}{@{}c r r r r l@{}}
\toprule
nome & \(N\) & \(A\) & \(B\) & \(C\) & \(\alpha\)\\
\midrule
+&2&1&6&216&\(3\sqrt3\)\\
+&4&2&15&1458&\(27/4\)\\
+&5&4&33&3375&\(15\sqrt3/2\)\\
-&9&1&5&-192&\(4\sqrt3/3\)\\
-&17&7&51&-1728&\(12\sqrt3\)\\
-&25&1&9&-8640&\(4\sqrt{15}/5\)\\
-&41&53&615&-110592&\(96\sqrt3\)\\
-&49&13&165&-326592&\(108\sqrt7/7\)\\
-&89&827&14151&-27000000&\(1500\sqrt3\)\\
\bottomrule
\end{tabular}
\end{table}

\begin{table}[htbp]
\centering
\caption{All primitive integer solutions for Level \(4A\) (taking \(B>0\))}\label{tab:4Acomplete}
\footnotesize
\setlength{\tabcolsep}{5pt}
\renewcommand{\arraystretch}{1.14}
\begin{tabular}{@{}c r r r r l l@{}}
\toprule
nome & \(N\) & \(A\) & \(B\) & \(C\) & \(\alpha\) & Convergence\\
\midrule
+&3&1&6&256&\(4\)&absolute\\
+&7&5&42&4096&\(16\)&absolute\\
-&2&1&4&-64&\(2\)&conditional\\
-&4&1&6&-512&\(2\sqrt2\)&absolute\\
\bottomrule
\end{tabular}
\end{table}
%must keep tags :(
%copy me: \tilde{\mathfrak F}
\subsection*{Euler-companion branch}
In this subsection we fix the meaning of ``companion'' to avoid the ambiguity caused by the not entirely consistent notation ``B'' and ``C'' branches in the literature. Set
\[
(K_1,K_2,K_3,K_4)=(432,64,27,16)=\frac14
(\kappa_1,\kappa_2,\kappa_3,\kappa_4).
\tag{1}
\]
For \(\ell=1,2,3,4\), define the Euler--companion coefficients
\[
U_\ell(n)=
\sum_{k=0}^{n}T_\ell(k)
\binom{n+k}{n-k}(-K_\ell)^{n-k}.
\tag{2}
\]
In this paper we denote these four sequences respectively by $1B,2B,3B,4C$. 
\begin{remark}
    In the notation of Chen--Glebov--Goenka, the last branch is called \(2C\); we use \(4C\) to remain consistent with the level \(4A\) single-sum sequence that generates it.
\end{remark}

\begin{lemma}[Euler generating function transform]\label{lem:eulertransform}
Let
\[
\widetilde {\mathfrak F}_\ell(Y)=\sum_{n\ge0}U_\ell(n)Y^n.
\]
Then in a neighbourhood of \(Y=0\) we have identically
\[
\widetilde {\mathfrak F}_\ell(Y)=\frac{1}{1+K_\ell Y}
{\mathfrak F}_\ell\!\left(\frac{Y}{(1+K_\ell Y)^2}\right).
\tag{3}
\]
\end{lemma}

\begin{proof}
From (2), put \(n=k+j\) and interchange summation gives
\[
\begin{aligned}
\widetilde {\mathfrak F}_\ell(Y)
&=\sum_{k\ge0}T_\ell(k)Y^k
  \sum_{j\ge0}\binom{2k+j}{j}(-K_\ell Y)^j\\
&=\sum_{k\ge0}T_\ell(k)Y^k(1+K_\ell Y)^{-2k-1}\\
&=\frac{1}{1+K_\ell Y}
  \sum_{k\ge0}T_\ell(k)
  \left(\frac{Y}{(1+K_\ell Y)^2}\right)^k.
\end{aligned}
\]
The last sum is precisely \({\mathfrak F}_\ell\).
\end{proof}

\begin{lemma}[Explicit transformations]\label{lem:companionslope}
Let
\[
Y=\frac1c,\qquad c\in\mathbf Z,\qquad c\ne-K_\ell,
\]
and define
\[
x=\frac{Y}{(1+K_\ell Y)^2}=\frac1J,
\qquad
{J=\frac{(c+K_\ell)^2}{c}}.
\tag{4}
\]
If
\[
\widetilde {\mathfrak S}_\ell(A_c,B_c,c)
=\sum_{n\ge0}U_\ell(n)\frac{A_c+B_cn}{c^n},
\]
then
\[
\widetilde {\mathfrak S}_\ell(A_c,B_c,c)
=\frac1J\bigl(a{\mathfrak F}_\ell(J^{-1})+b\Theta {\mathfrak F}_\ell(J^{-1})\bigr),
\tag{5}
\]
where
\[
a=A_c(c+K_\ell)-B_cK_\ell,
\qquad
b=B_c(c-K_\ell).
\tag{6}
\]
If \(B_c\ne0\) and \(c\ne K_\ell\), write
\(\lambda_c=A_c/B_c\) and \(\lambda_A=a/b\). Then
\[
{
\lambda_A=
\frac{(c+K_\ell)\lambda_c-K_\ell}{c-K_\ell},
\qquad
\lambda_c=
\frac{K_\ell+(c-K_\ell)\lambda_A}{c+K_\ell}.
}
\tag{7}
\]
Let $g=\gcd(a,b)>0$ and write $a=gA$, $b=gB$ with $\gcd(A,B)=1$. If
\[
A{\mathfrak F}_\ell(J^{-1})+B\Theta {\mathfrak F}_\ell(J^{-1})=\frac{\alpha_A}{\pi},
\]
then
\[
{\alpha_c=\frac{g}{J}\alpha_A.}
\tag{8}
\]
\end{lemma}

\begin{proof}
Apply \(\Theta_Y=Y\,d/dY\) to (3). If
\(r=1+K_\ell Y\), then
\[
Y\frac{d}{dY}\log\frac{Y}{r^2}
=\frac{1-K_\ell Y}{1+K_\ell Y}.
\]
Therefore
\[
\Theta_Y\widetilde {\mathfrak F}_\ell(Y)
=-\frac{K_\ell Y}{r^2}{\mathfrak F}_\ell(x)
+\frac{1-K_\ell Y}{r^2}\Theta {\mathfrak F}_\ell(x).
\]
Substituting \(Y=1/c\) into
\(A_c\widetilde {\mathfrak F}_\ell+B_c\Theta_Y\widetilde {\mathfrak F}_\ell\) and simplifying gives (5)--(6). Equation (7) is the simplification of $a/b$ and its inverse.  After writing $a=gA$ and $b=gB$ with $g>0$ and $\gcd(A,B)=1$, equation (5) immediately gives (8).
\end{proof}

\begin{corollary}[Discriminant condition for integer denominators]\label{cor:disccomp}
If \(J\in\mathbb Z\), then (4) is equivalent to
\[
c^2+(2K_\ell-J)c+K_\ell^2=0.
\tag{9}
\]
Hence a necessary condition for an integer \(c\) to exist is
\[
{\Delta=J(J-4K_\ell)\ \text{is an integer square}.}
\tag{10}
\]
Conversely, if \(\Delta\) is a square and
\((J-2K_\ell\pm\sqrt\Delta)/2\in\mathbb Z\), then the two integer roots are
\[
c=\frac{J-2K_\ell\pm\sqrt{J(J-4K_\ell)}}{2}.
\tag{11}
\]
\end{corollary}

%\begin{remark}\label{rem:Jnotautomatic}
Corollary~\ref{cor:disccomp} is only an integrality filter for those cases in which $J\in\mathbb Z$. Completeness below is instead based on the direct CM enumeration of the companion uniformizer values.
%\end{remark}

\begin{lemma}[Convergence boundary for the companion]\label{lem:compconv}
For \(U_\ell(n)\) there exists a constant \(d_\ell>0\) such that the leading asymptotics have sign
\((-1)^n\), and
\[
U_\ell(n)=(-1)^nK_\ell^n\times
\begin{cases}
 d_1n^{-1/3}(1+o(1)),&\ell=1,\\
 d_2n^{-1/2}(1+o(1)),&\ell=2,\\
 d_3n^{-2/3}(1+o(1)),&\ell=3,\\
 d_4\dfrac{\log n}{n}(1+o(1)),&\ell=4.
\end{cases}
\tag{12}
\]
Therefore
\[
\begin{array}{c|c}
|c|>K_\ell&\text{absolutely convergent},\\
 c=K_\ell&\text{convergent iff }B_c=0\text{, and only conditionally},\\
 c=-K_\ell&\text{divergent},\\
 |c|<K_\ell&\text{divergent}.
\end{array}
\tag{13}
\]
\end{lemma}

\begin{proof}
The identity
\[
1-\kappa_\ell\frac{Y}{(1+K_\ell Y)^2}
=\left(\frac{1-K_\ell Y}{1+K_\ell Y}\right)^2
\]
shows why the other point on the same circle, $Y=K_\ell^{-1}$, is regular after the quadratic pullback: the local square-root branch at $\kappa_\ell x=1$ becomes an ordinary local parameter.  Thus the genuine dominant singularity is $Y=-K_\ell^{-1}$.  Put $r=1+K_\ell Y$. As $r\to0$,
\[
\frac{Y}{r^2}\longrightarrow-\infty.
\]
Using the standard connection formula for
\[
{}_3F_2\left(a_\ell,\frac12,1-a_\ell;1,1;\kappa_\ell x\right)
\]
as \(x\to-\infty\): when
\(a_\ell=1/6,1/4,1/3\), the minimal parameter term gives respectively singularities of order
\(r^{-2/3},r^{-1/2},r^{-1/3}\); when \(a_4=1/2\) there is a triple coincidence,
and the dominant singularity is \((\log r)^2\). The Darboux--Flajolet transfer theorem then gives (12).

When \(|c|>K_\ell\), the root test gives absolute convergence; when \(|c|<K_\ell\), the general term does not tend to zero.
The singularity analysis can be sharpened to a full asymptotic expansion in descending powers of $n$ (with powers of $\log n$ in the repeated-parameter case), so the normalized absolute values are eventually decreasing.  At $c=K_\ell$, the factor $(-1)^n$ therefore gives convergence of the $A_c$-part by the alternating test, whereas if $B_c\ne0$, multiplication by $n$ prevents the general term from tending to zero.
At \(c=-K_\ell\), the alternating sign is cancelled; the \(A_c\)-part behaves respectively as
\(n^{-1/3},n^{-1/2},n^{-2/3},(\log n)/n\), all divergent;
if \(B_c\ne0\) convergence is even more impossible.
\end{proof}

\begin{theorem}[Complete primitive integer classification of Euler--companion branches]\label{thm:compcomplete}
Under \((\mathrm{GPC}_{\mathrm{tor}})\), in the Euler--companion branches \(1B,2B,3B,4C\) as defined in (2), up to the overall sign there are exactly
\(5\) convergent primitive integer solutions, all listed in Table~\ref{tab:companion}.
Among them \(1B\) has no solution, \(2B\) has three, \(3B\) and \(4C\) each have one; all five are absolutely convergent.
\end{theorem}

\begin{proof}
First, by Lemma~\ref{lem:compconv}, \(|c|<K_\ell\) and \(c=-K_\ell\) are all excluded.
If \(c=K_\ell\), convergence forces \(B_c=0\). From (3), then
\[
\widetilde {\mathfrak F}_\ell(K_\ell^{-1})=\frac12{\mathfrak F}_\ell(\kappa_\ell^{-1}),
\]
so if it were an algebraic number times \(1/\pi\), this would contradict Lemma~\ref{lem:Apositiveboundary}.
Thus we only need consider \(|c|>K_\ell\).

For such a solution, (4) gives a rational A-side modular parameter $J$, and (5) converts the companion relation into an A-side linear period relation in the same Picard--Fuchs system; hence the period-torsor hypothesis again forces the corresponding point to be CM.  To obtain a complete list of integer companion denominators, we use the direct rational CM uniformizer enumeration in Chen--Glebov--Goenka: their Table~15 gives the rational $2B,3B$ Euler cases, and their Table~13 together with equation (49) gives the rational $2C$ (our $4C$) case~\cite[Tables~13,15]{CGG}.  This direct enumeration is essential in view of Remark~\ref{rem:Jnotautomatic}.  Translating those companion parameters through (4), and excluding the boundary values
\(J=0\) (corresponding to \(c=-K_\ell\)) and \(J=4K_\ell\) (corresponding to \(c=K_\ell\)),
the non-boundary CM parameters that can produce integer \(c\) are exactly
\[
\begin{array}{c|c|c}
\text{Branch}&J&\text{Integer roots }c\text{ from (11)}\\ \hline
1B&\text{none}&\text{none}\\
2B&648&8,\ 512\\
2B&-144&-16,\ -256\\
2B&-3969&-1,\ -4096\\
3B&-192&-3,\ -243\\
4C&-8&-8,\ -32.
\end{array}
\tag{14}
\]
Lemma~\ref{lem:compconv} retains only
\[
512,\\ -256,\\ -4096,\\ -243,\\ -32.
\tag{15}
\]
The corresponding primitive slopes and right-hand constants are obtained from the same two tables of Chen--Glebov--Goenka;
alternatively, in the first three and the fourth cases they may be checked consistently from the corresponding A-side slopes using Lemma~\ref{lem:companionslope}. This yields the five rows of Table~\ref{tab:companion}.
Since each row satisfies \(|c|>K_\ell\), all five series are absolutely convergent.
The complete CM enumeration together with the one-dimensionality of the relation space for fixed parameters excludes any other primitive slopes.
\end{proof}

\begin{table}[htbp]
\centering
\caption{All Euler--companion primitive integer solutions (taking \(B>0\))}\label{tab:companion}
\small
\setlength{\tabcolsep}{5pt}
\renewcommand{\arraystretch}{1.18}
\begin{tabular}{@{}c r r r l l@{}}
\toprule
Branch & \(A\) & \(B\) & \(C\) & \(\alpha\) & Reference\\
\midrule
\(1B\)&\multicolumn{4}{c}{no convergent primitive integer solution}&\\
\(2B\)&2&9&512&\(4\)&\cite[Table~15]{CGG}\\
\(2B\)&0&1&-256&\(4\sqrt3/9\)&\cite[Table~15]{CGG}\\
\(2B\)&1&9&-4096&\(64\sqrt7/49\)&\cite[Table~15]{CGG}\\
\(3B\)&1&8&-243&\(3\sqrt3\)&\cite[Table~15]{CGG}\\
\(4C\)&0&1&-32&\(2\)&\cite[Table~13, Eq.~(49)]{CGG}\\
\bottomrule
\end{tabular}
\end{table}

The corresponding five identities can be written uniformly as
\[
\sum_{n=0}^{\infty}U_\ell(n)\frac{A+Bn}{C^n}=\frac{\alpha}{\pi}.
\tag{16}
\]
In particular, level \(4C\) can also be written in a pure convolution form. Indeed,
\[
U_4(n)=(-1)^n
\sum_{k=0}^{n}
\binom{2k}{k}^{2}
\binom{2n-2k}{n-k}^{2},
\tag{17}
\]
so the last row is equivalent to
\[
{
\sum_{n=0}^{\infty}\frac{n}{32^n}
\sum_{k=0}^{n}
\binom{2k}{k}^{2}
\binom{2n-2k}{n-k}^{2}
=\frac2\pi.
}
\tag{18}
\]

\begin{corollary}[Total classification for levels \(1\)--\(4\)]\label{cor:41}
Under the GPC assumption and the branch definitions in this paper, the convergent primitive integer \(1/\pi\) series in the standard A branches and the Euler--companion branches have, up to the overall sign, exactly
\[
{36+5=41}
\]
sign classes, with branch counts
\[
\begin{array}{c|cccccccc}
\text{Branch}&1A&1B&2A&2B&3A&3B&4A&4C\\ \hline
\text{Count}&11&0&12&3&9&1&4&1.
\end{array}
\tag{19}
\]
If \((A,B,C)\) and \((-A,-B,C)\) are regarded as distinct integer triples, then there are \(82\) triples.
For all \(41\) representatives, \(\alpha\) satisfies
\([\mathbf Q(\alpha):\mathbf Q]\le2\), but this degree bound is not an assumption of the GPC finiteness argument.
\end{corollary}

%\section{References}

\begin{thebibliography}{99}

\bibitem{BCC}
J.-B. Bost and F. Charles,
\emph{Some remarks concerning the Grothendieck period conjecture},
J. Reine Angew. Math. 714 (2016), 175--208;
arXiv:1307.1045.

\bibitem{CCY}
J. M. Campbell, S. Cooper and D. Ye,
\emph{Quadratic irrational analogues of Ramanujan's series for \(1/\pi\)},
arXiv:2602.09352 (2026).

\bibitem{CGG}
I. Chen, G. Glebov and R. Goenka,
\emph{Chudnovsky--Ramanujan type formulae for non-compact arithmetic triangle groups},
J. Number Theory 241 (2022), 603--654;
arXiv:1909.07350.

\bibitem{Andre}
Y. Andr\'e,
\emph{Une introduction aux motifs},
Panoramas et Synth\`eses 17, Soci\'et\'e Math\'ematique de France, 2004.

\bibitem{Cox}
D. A. Cox,
\emph{Primes of the Form \(x^2+ny^2\)},
2nd ed., Wiley, 2013.

\bibitem{Katz}
N. Katz,
\emph{\(p\)-adic properties of modular schemes and modular forms},
in Modular Functions of One Variable III,
Lecture Notes in Mathematics 350, Springer, 1973.

\bibitem{Ram}
S.~Ramanujan, ``Modular equations and approximations to $\pi$,''
\emph{Quart. J. Math. (Oxford)}, vol.~45, pp.~350--372, 1914.
Reprinted in [93, 23--39].

\end{thebibliography}

\end{document}
